\section{Introduction}

Let \(C=\{\sum_{j\ge1} a_j3^{-j}:a_j\in\{0,2\}\}\) and
\(\kappa=\log 2/\log 3\). For \(T\ge2\), put
\[
 N_C(T)=\#\{a/q\in C:(a,q)=1,\ 1\le q\le T\}.
\]
We count distinct fractions, including \(0\) and \(1\). The original
symbolic denominator of a periodic digit expansion can substantially
exceed the reduced denominator. Counting short expansions alone therefore
does not give the required upper bound.

\begin{theorem}\label{thm:main}
For every \(\varepsilon>0\) one has
\[
                 N_C(T)\ll_\varepsilon T^{\kappa+\varepsilon}
                 \qquad(T\ge2).
\]
\end{theorem}

The theorem proves Conjecture~3.3 of Broderick, Fishman, and
Reich~\cite{BFR2011}. Their formulation uses the shells
\(3^j\le q\le3^{j+1}\): the cumulative and shell assertions follow
from each other by restriction and geometric summation, respectively,
with the positive exponent slack rescaled by \(\kappa\). It is the
middle-third case of \cite[Conjecture~1.1]{CVY2026}, not an assertion
for other missing-digit sets. The conjecture grew out of
Mahler's question about rational approximation inside \(C\)~\cite{Mahler1984}.
Every two distinct fractions up to denominator \(T\) are \(T^{-2}\)-separated;
ordinary packing in \(C\) thus gives \(O(T^{2\kappa})\). Theorem~\ref{thm:main}
identifies the power at the scale of the denominator itself.

Rahm, Solomon, Trauthwein, and Weiss proved
\(N_C(T)\gg T^\kappa\log T\)~\cite[Proposition~3.3]{RSTW2020}.
One can see the origin of the logarithm directly. Set \(L=\lfloor\log_3 T\rfloor\).
For \(1\le i<L\), take the expansions \(0.A\overline B\) with prefix
\(A\) of length \(i\), period \(B\) of length \(L-i\), and different final
trits in the two words. The numerator before reduction is
\(A(3^{L-i}-1)+B\), viewing the words as base-three integers. It is
prime to three, so the reduced denominator has exact \(3\)-adic
valuation \(i\), and is at most \(3^L\). For each \(i\) there are
\(2^{L-1}\) points, since different all-\(\{0,2\}\) addresses represent
distinct points. Thus the power of \(T\) in the theorem is sharp,
but the \(\varepsilon\) cannot simply be deleted.

\subsection*{The local counting problem}

First consider denominators \(q\) prime to \(3\). The action of
multiplication by \(3\) on \(a/q\) has order \(n=\operatorname{ord}_q(3)\),
and every iterate stays in \(C\). We descend through divisors \(h\) of
\(n\), retaining the real difference \(\{3^h x\}-x\), where braces
denote fractional part. Its denominator is \(q/(q,3^h-1)\).
If a descent step increases this denominator by a factor \(u\),
then the child differences at fixed \(u\), with the preceding real
difference fixed, sit on a single translate of the order-\(u\) circle
grid. The product of the factors \(u\) and the remaining terminal
factor is \(q\).

Here is the quantitative balance behind the proof. At \(u\asymp U\)
the difference set of the child grid is enlarged to a small-doubling,
tripling-invariant set of size \(U^\beta\), say. Boolean sumset growth
bounds the number of child values by \(U^{\kappa\beta+o(1)}\), with
arbitrarily small power slack. To count the moduli \(u\) attaining that
size, split each into entire prime-power factors. Iterated addition
and rotation show that the factors with substantial digit entropy have
almost full combined projection. A common detector on least common
multiples lets us count those nearly full projections at a fixed
complement. The number of names for the complementary factors then
costs at most \(U^{(1-\beta)/2}\), up to small losses, because an inactive
prime leaves a deficit of roughly the logarithm of that prime. Hence
\[
                  \kappa\beta+\tfrac12(1-\beta)\le\kappa .
\]
The bound is for the sum over moduli and child values at one actual
parent, not a uniform assertion about a single denominator.

Section~\ref{sec:itinerary} constructs the descent and accounts for the
short steps jointly. Sections~\ref{sec:detector} and~\ref{sec:native}
give the common detector, which works even when the independently chosen
grid fibers at a common lag have different offsets. The Boolean estimate
and its least-common-multiple application are in
Sections~\ref{sec:boolean}--\ref{sec:lcm}. Section~\ref{sec:tariff}
pays for the complementary primes, and Section~\ref{sec:assembly}
assembles the counting argument and restores preperiods.

We write \(G_u=u^{-1}\mathbb Z/\mathbb Z\) and use additive notation on
the circle. If \(A\) is a set and \(m\) a nonnegative integer, \(mA\)
denotes its \(m\)-fold sumset while \(m\cdot A=\{ma:a\in A\}\)
denotes its scalar image. All entropies and unmarked logarithms use
natural logarithms. Implied constants may depend on fixed tolerances,
but not on the height or the choices of affine offsets.
